\documentclass[11pt,a4paper]{article}
\usepackage[utf8]{inputenc}
\usepackage[T1]{fontenc}
\usepackage[catalan]{babel}
\usepackage{lmodern}
\usepackage[sfdefault,lining]{FiraSans}
\usepackage[a4paper,top=2.3cm,bottom=2.2cm,left=1.8cm,right=1.8cm,headheight=1.3cm,headsep=0.4cm,footskip=1cm]{geometry}
\usepackage{amsmath,amssymb}
\usepackage{icomma}
\usepackage{graphicx}
\usepackage[table]{xcolor}
\usepackage{tikz}
\usepackage{fontawesome5}
\usepackage{ccicons}
\usepackage{fancyhdr}
\usepackage{lastpage}
\usepackage{tasks}
\usepackage{needspace}
\usepackage[most]{tcolorbox}
\usepackage{array}
\usepackage[hidelinks,pdftitle={Fitxa d'exercicis de la unitat 5: Proporcionalitat i percentatges. Nombres decimals. 1r d'ESO},pdfsubject={Matemàtiques 1r ESO},pdfkeywords={nombres decimals, aproximació, operacions amb decimals, decimals i fraccions, raons, proporcions, proporcionalitat directa, proporcionalitat inversa, regla de tres, percentatges}]{hyperref}

\definecolor{navy}{HTML}{1B2748}
\definecolor{accent}{HTML}{E8603A}
\definecolor{teal}{HTML}{1F9E90}
\definecolor{soft}{HTML}{EEF1F8}
\definecolor{muted}{HTML}{5D6685}

\setlength{\parindent}{0pt}
\setlength{\parskip}{3pt}
\settasks{label=\alph*), label-format=\bfseries\color{navy}, label-width=1.3em, item-indent=1.8em, column-sep=0.6em, after-item-skip=5pt, before-skip=3pt, after-skip=4pt}

% ---------- capçalera i peu ----------
\newcommand{\llicencia}{\href{https://creativecommons.org/licenses/by-nc-sa/4.0/deed.ca}{\ccbyncsa}\ \ Llicència Creative Commons Reconeixement-NoComercial-CompartirIgual 4.0 Internacional (CC BY-NC-SA 4.0)}
\pagestyle{fancy}
\fancyhf{}
\fancyhead[L]{\raisebox{-0.35\height}{\includegraphics[height=0.95cm]{logo-gva.png}}}
\fancyhead[C]{\small\color{muted}Fitxa d'exercicis · Unitat 5 · 1r d'ESO}
\fancyhead[R]{\raisebox{-0.35\height}{\includegraphics[height=0.95cm]{logo-vermellar.png}}}
\fancyfoot[L]{\scriptsize\color{muted}\llicencia}
\fancyfoot[R]{\small\color{muted}\thepage}
\renewcommand{\headrulewidth}{0.4pt}
\renewcommand{\headrule}{\color{muted!40}\hrule height \headrulewidth}
\fancypagestyle{primera}{\fancyhf{}\renewcommand{\headrulewidth}{0pt}
  \fancyfoot[L]{\scriptsize\color{muted}\llicencia}
  \fancyfoot[R]{\small\color{muted}\thepage}}

% ---------- nivells i exercicis ----------
\newcommand{\estrelles}[1]{{\color{accent}\ifcase#1\or\faStar\faStar[regular]\faStar[regular]\or\faStar\faStar\faStar[regular]\or\faStar\faStar\faStar\fi}}
\newcounter{exer}
\newcommand{\ex}[1]{\par\needspace{4\baselineskip}\medskip\stepcounter{exer}\noindent{\color{navy}\bfseries\theexer.}\ \ {\small\estrelles{#1}}\ \ }

\newcommand{\seccio}[2]{\par\needspace{10\baselineskip}\bigskip
  \begin{tcolorbox}[colback=navy,colframe=navy,coltext=white,arc=2mm,boxrule=0pt,left=3mm,right=3mm,top=1.2mm,bottom=1.2mm]
  {\large\bfseries #1\quad #2}
  \end{tcolorbox}\vspace{-2pt}}

\newcolumntype{C}[1]{>{\centering\arraybackslash}p{#1}}

% quadre de recordatori (només a la fitxa d'exercicis)
\newcommand{\recorda}[1]{\par\smallskip\begin{tcolorbox}[colback=accent!7,colframe=accent!55,arc=2mm,boxrule=0.5pt,left=2.5mm,right=2.5mm,top=1mm,bottom=1mm]\small{\color{accent}\faLightbulb[regular]}\ \textbf{Recorda:}\ #1\end{tcolorbox}}
\newcommand{\res}[1]{#1}\newcommand{\rest}[1]{#1}

% =====================================================================
% Capçalera FentMates de les fitxes: un color per a cada curs.
% Els colors són els de la pàgina inicial (LEVELS a assets/js/materials.js);
% el logo ve de _eines/logo/pdf/ (python3 _eines/logo/genera.py).
%
% Ús, al final del preàmbul de la unitat:
%   % =====================================================================
% Capçalera FentMates de les fitxes: un color per a cada curs.
% Els colors són els de la pàgina inicial (LEVELS a assets/js/materials.js);
% el logo ve de _eines/logo/pdf/ (python3 _eines/logo/genera.py).
%
% Ús, al final del preàmbul de la unitat:
%   % =====================================================================
% Capçalera FentMates de les fitxes: un color per a cada curs.
% Els colors són els de la pàgina inicial (LEVELS a assets/js/materials.js);
% el logo ve de _eines/logo/pdf/ (python3 _eines/logo/genera.py).
%
% Ús, al final del preàmbul de la unitat:
%   \input{../comu/fentmates.tex}
%   \fmcurs{1eso}
% i, al començament del document:
%   \capcalera{Matemàtiques}{1r ESO}{Fitxa d'exercicis · Unitat 2}{Divisibilitat}
% Opcional abans de \input: \def\fmneed{14} (espai mínim abans de cada secció).
% =====================================================================
\makeatletter
\def\fm@logodir{../../logo/pdf}
\providecommand{\fmneed}{10}
\def\fm@cursnom{}

% clau del curs, color i color del text damunt del color
\newcommand{\fm@curs}[3]{\@namedef{fm@color@#1}{#2}\@namedef{fm@text@#1}{#3}}
\fm@curs{1eso}{E6B84A}{navy}
\fm@curs{2eso}{3B82F6}{white}
\fm@curs{3eso}{FF7A50}{white}
\fm@curs{4eso-a}{A78BFA}{white}
\fm@curs{4eso-b}{6D28D9}{white}
\fm@curs{1batx-mi}{9B1C1C}{white}
\fm@curs{1batx-ccss}{EF6B6B}{white}
\fm@curs{2batx-mii}{1E6B47}{white}
\fm@curs{2batx-ccss}{2BB3A3}{white}

\newcommand{\fmcurs}[1]{%
  \@ifundefined{fm@color@#1}{\PackageError{fentmates}{Curs desconegut: #1}{}}{}%
  \edef\fm@tmp{\noexpand\definecolor{curs}{HTML}{\@nameuse{fm@color@#1}}%
               \noexpand\colorlet{curstext}{\@nameuse{fm@text@#1}}}\fm@tmp
  \colorlet{cursfosc}{curs!72!black}%
  \def\fm@logo{\fm@logodir/vertical-#1.pdf}%
  \def\fm@marca{\fm@logodir/marca-#1.pdf}}

% ---------- capçalera de la primera pàgina ----------
% #1 matèria (petit), #2 curs, #3 tipus de document i unitat, #4 títol
\newcommand{\capcalera}[4]{%
  \gdef\fm@cursnom{#2}%
  \par\noindent
  \begin{tcolorbox}[colback=curs,colframe=curs,boxrule=0pt,arc=2.2mm,width=3.5cm,height=2.15cm,
      valign=center,left=3mm,right=2mm,top=1mm,bottom=1mm,nobeforeafter,box align=center]
    \color{curstext}{\footnotesize #1}\par\vspace{1pt}{\fontsize{19}{21}\selectfont\bfseries #2\par}%
  \end{tcolorbox}\hspace{5mm}%
  \begin{minipage}[c]{\dimexpr\linewidth-3.5cm-5mm-2.9cm\relax}
    {\color{cursfosc}\small\bfseries\MakeUppercase{#3}}\par\vspace{2pt}
    {\color{navy}\fontsize{20}{23}\selectfont\bfseries #4\par}
  \end{minipage}\hfill
  \begin{minipage}[c]{2.6cm}\raggedleft\includegraphics[width=2.6cm]{\fm@logo}\end{minipage}\par
  \vspace{2pt}}

% ---------- barres de secció del color del curs ----------
\renewcommand{\seccio}[2]{\par\needspace{\fmneed\baselineskip}\bigskip
  \begin{tcolorbox}[colback=curs,colframe=curs,coltext=curstext,arc=2mm,boxrule=0pt,left=3mm,right=3mm,top=1.2mm,bottom=1.2mm]
  {\large\bfseries #1\quad #2}
  \end{tcolorbox}\vspace{-2pt}}

% ---------- estrelles de dificultat del color del curs (a Batxillerat no se'n posen) ----------
\@ifundefined{estrelles}{}{%
  \renewcommand{\estrelles}[1]{{\color{curs}\ifcase#1\or\faStar\faStar[regular]\faStar[regular]\or\faStar\faStar\faStar[regular]\or\faStar\faStar\faStar\fi}}}

% ---------- capçalera i peu de les pàgines ----------
% primera pàgina: logos institucionals al peu
\fancypagestyle{primera}{\fancyhf{}\renewcommand{\headrulewidth}{0pt}
  \fancyfoot[L]{\raisebox{-0.32\height}{\includegraphics[height=0.78cm]{logo-gva.png}}\hspace{5mm}%
    \raisebox{-0.32\height}{\includegraphics[height=0.78cm]{logo-vermellar.png}}}
  \fancyfoot[R]{\scriptsize\color{muted}\href{https://creativecommons.org/licenses/by-nc-sa/4.0/deed.ca}{\ccbyncsa}\ \ CC BY-NC-SA 4.0 · fentmates.org\qquad\small\thepage}}
% la resta: el símbol a l'esquerra i el curs a la dreta, amb el títol al centre com abans
\fancyhead[L]{\raisebox{-0.3\height}{\includegraphics[height=0.62cm]{\fm@marca}}}
\fancyhead[R]{\small\bfseries\color{cursfosc}\fm@cursnom}
\renewcommand{\headrule}{\color{curs!45}\hrule height \headrulewidth}
\makeatother

%   \fmcurs{1eso}
% i, al començament del document:
%   \capcalera{Matemàtiques}{1r ESO}{Fitxa d'exercicis · Unitat 2}{Divisibilitat}
% Opcional abans de \input: \def\fmneed{14} (espai mínim abans de cada secció).
% =====================================================================
\makeatletter
\def\fm@logodir{../../logo/pdf}
\providecommand{\fmneed}{10}
\def\fm@cursnom{}

% clau del curs, color i color del text damunt del color
\newcommand{\fm@curs}[3]{\@namedef{fm@color@#1}{#2}\@namedef{fm@text@#1}{#3}}
\fm@curs{1eso}{E6B84A}{navy}
\fm@curs{2eso}{3B82F6}{white}
\fm@curs{3eso}{FF7A50}{white}
\fm@curs{4eso-a}{A78BFA}{white}
\fm@curs{4eso-b}{6D28D9}{white}
\fm@curs{1batx-mi}{9B1C1C}{white}
\fm@curs{1batx-ccss}{EF6B6B}{white}
\fm@curs{2batx-mii}{1E6B47}{white}
\fm@curs{2batx-ccss}{2BB3A3}{white}

\newcommand{\fmcurs}[1]{%
  \@ifundefined{fm@color@#1}{\PackageError{fentmates}{Curs desconegut: #1}{}}{}%
  \edef\fm@tmp{\noexpand\definecolor{curs}{HTML}{\@nameuse{fm@color@#1}}%
               \noexpand\colorlet{curstext}{\@nameuse{fm@text@#1}}}\fm@tmp
  \colorlet{cursfosc}{curs!72!black}%
  \def\fm@logo{\fm@logodir/vertical-#1.pdf}%
  \def\fm@marca{\fm@logodir/marca-#1.pdf}}

% ---------- capçalera de la primera pàgina ----------
% #1 matèria (petit), #2 curs, #3 tipus de document i unitat, #4 títol
\newcommand{\capcalera}[4]{%
  \gdef\fm@cursnom{#2}%
  \par\noindent
  \begin{tcolorbox}[colback=curs,colframe=curs,boxrule=0pt,arc=2.2mm,width=3.5cm,height=2.15cm,
      valign=center,left=3mm,right=2mm,top=1mm,bottom=1mm,nobeforeafter,box align=center]
    \color{curstext}{\footnotesize #1}\par\vspace{1pt}{\fontsize{19}{21}\selectfont\bfseries #2\par}%
  \end{tcolorbox}\hspace{5mm}%
  \begin{minipage}[c]{\dimexpr\linewidth-3.5cm-5mm-2.9cm\relax}
    {\color{cursfosc}\small\bfseries\MakeUppercase{#3}}\par\vspace{2pt}
    {\color{navy}\fontsize{20}{23}\selectfont\bfseries #4\par}
  \end{minipage}\hfill
  \begin{minipage}[c]{2.6cm}\raggedleft\includegraphics[width=2.6cm]{\fm@logo}\end{minipage}\par
  \vspace{2pt}}

% ---------- barres de secció del color del curs ----------
\renewcommand{\seccio}[2]{\par\needspace{\fmneed\baselineskip}\bigskip
  \begin{tcolorbox}[colback=curs,colframe=curs,coltext=curstext,arc=2mm,boxrule=0pt,left=3mm,right=3mm,top=1.2mm,bottom=1.2mm]
  {\large\bfseries #1\quad #2}
  \end{tcolorbox}\vspace{-2pt}}

% ---------- estrelles de dificultat del color del curs (a Batxillerat no se'n posen) ----------
\@ifundefined{estrelles}{}{%
  \renewcommand{\estrelles}[1]{{\color{curs}\ifcase#1\or\faStar\faStar[regular]\faStar[regular]\or\faStar\faStar\faStar[regular]\or\faStar\faStar\faStar\fi}}}

% ---------- capçalera i peu de les pàgines ----------
% primera pàgina: logos institucionals al peu
\fancypagestyle{primera}{\fancyhf{}\renewcommand{\headrulewidth}{0pt}
  \fancyfoot[L]{\raisebox{-0.32\height}{\includegraphics[height=0.78cm]{logo-gva.png}}\hspace{5mm}%
    \raisebox{-0.32\height}{\includegraphics[height=0.78cm]{logo-vermellar.png}}}
  \fancyfoot[R]{\scriptsize\color{muted}\href{https://creativecommons.org/licenses/by-nc-sa/4.0/deed.ca}{\ccbyncsa}\ \ CC BY-NC-SA 4.0 · fentmates.org\qquad\small\thepage}}
% la resta: el símbol a l'esquerra i el curs a la dreta, amb el títol al centre com abans
\fancyhead[L]{\raisebox{-0.3\height}{\includegraphics[height=0.62cm]{\fm@marca}}}
\fancyhead[R]{\small\bfseries\color{cursfosc}\fm@cursnom}
\renewcommand{\headrule}{\color{curs!45}\hrule height \headrulewidth}
\makeatother

%   \fmcurs{1eso}
% i, al començament del document:
%   \capcalera{Matemàtiques}{1r ESO}{Fitxa d'exercicis · Unitat 2}{Divisibilitat}
% Opcional abans de \input: \def\fmneed{14} (espai mínim abans de cada secció).
% =====================================================================
\makeatletter
\def\fm@logodir{../../logo/pdf}
\providecommand{\fmneed}{10}
\def\fm@cursnom{}

% clau del curs, color i color del text damunt del color
\newcommand{\fm@curs}[3]{\@namedef{fm@color@#1}{#2}\@namedef{fm@text@#1}{#3}}
\fm@curs{1eso}{E6B84A}{navy}
\fm@curs{2eso}{3B82F6}{white}
\fm@curs{3eso}{FF7A50}{white}
\fm@curs{4eso-a}{A78BFA}{white}
\fm@curs{4eso-b}{6D28D9}{white}
\fm@curs{1batx-mi}{9B1C1C}{white}
\fm@curs{1batx-ccss}{EF6B6B}{white}
\fm@curs{2batx-mii}{1E6B47}{white}
\fm@curs{2batx-ccss}{2BB3A3}{white}

\newcommand{\fmcurs}[1]{%
  \@ifundefined{fm@color@#1}{\PackageError{fentmates}{Curs desconegut: #1}{}}{}%
  \edef\fm@tmp{\noexpand\definecolor{curs}{HTML}{\@nameuse{fm@color@#1}}%
               \noexpand\colorlet{curstext}{\@nameuse{fm@text@#1}}}\fm@tmp
  \colorlet{cursfosc}{curs!72!black}%
  \def\fm@logo{\fm@logodir/vertical-#1.pdf}%
  \def\fm@marca{\fm@logodir/marca-#1.pdf}}

% ---------- capçalera de la primera pàgina ----------
% #1 matèria (petit), #2 curs, #3 tipus de document i unitat, #4 títol
\newcommand{\capcalera}[4]{%
  \gdef\fm@cursnom{#2}%
  \par\noindent
  \begin{tcolorbox}[colback=curs,colframe=curs,boxrule=0pt,arc=2.2mm,width=3.5cm,height=2.15cm,
      valign=center,left=3mm,right=2mm,top=1mm,bottom=1mm,nobeforeafter,box align=center]
    \color{curstext}{\footnotesize #1}\par\vspace{1pt}{\fontsize{19}{21}\selectfont\bfseries #2\par}%
  \end{tcolorbox}\hspace{5mm}%
  \begin{minipage}[c]{\dimexpr\linewidth-3.5cm-5mm-2.9cm\relax}
    {\color{cursfosc}\small\bfseries\MakeUppercase{#3}}\par\vspace{2pt}
    {\color{navy}\fontsize{20}{23}\selectfont\bfseries #4\par}
  \end{minipage}\hfill
  \begin{minipage}[c]{2.6cm}\raggedleft\includegraphics[width=2.6cm]{\fm@logo}\end{minipage}\par
  \vspace{2pt}}

% ---------- barres de secció del color del curs ----------
\renewcommand{\seccio}[2]{\par\needspace{\fmneed\baselineskip}\bigskip
  \begin{tcolorbox}[colback=curs,colframe=curs,coltext=curstext,arc=2mm,boxrule=0pt,left=3mm,right=3mm,top=1.2mm,bottom=1.2mm]
  {\large\bfseries #1\quad #2}
  \end{tcolorbox}\vspace{-2pt}}

% ---------- estrelles de dificultat del color del curs (a Batxillerat no se'n posen) ----------
\@ifundefined{estrelles}{}{%
  \renewcommand{\estrelles}[1]{{\color{curs}\ifcase#1\or\faStar\faStar[regular]\faStar[regular]\or\faStar\faStar\faStar[regular]\or\faStar\faStar\faStar\fi}}}

% ---------- capçalera i peu de les pàgines ----------
% primera pàgina: logos institucionals al peu
\fancypagestyle{primera}{\fancyhf{}\renewcommand{\headrulewidth}{0pt}
  \fancyfoot[L]{\raisebox{-0.32\height}{\includegraphics[height=0.78cm]{logo-gva.png}}\hspace{5mm}%
    \raisebox{-0.32\height}{\includegraphics[height=0.78cm]{logo-vermellar.png}}}
  \fancyfoot[R]{\scriptsize\color{muted}\href{https://creativecommons.org/licenses/by-nc-sa/4.0/deed.ca}{\ccbyncsa}\ \ CC BY-NC-SA 4.0 · fentmates.org\qquad\small\thepage}}
% la resta: el símbol a l'esquerra i el curs a la dreta, amb el títol al centre com abans
\fancyhead[L]{\raisebox{-0.3\height}{\includegraphics[height=0.62cm]{\fm@marca}}}
\fancyhead[R]{\small\bfseries\color{cursfosc}\fm@cursnom}
\renewcommand{\headrule}{\color{curs!45}\hrule height \headrulewidth}
\makeatother

\fmcurs{1eso}

\begin{document}
\thispagestyle{primera}

\capcalera{Matemàtiques}{1r ESO}{Fitxa d'exercicis · Unitat 5}{Proporcionalitat i percentatges. Nombres decimals}
\vspace{10pt}
Nom i cognoms: \rule{8.4cm}{0.4pt}\hfill Grup: \rule{1.6cm}{0.4pt}\hfill Data: \rule{2.2cm}{0.4pt}
\vspace{8pt}

\begin{tcolorbox}[colback=soft,colframe=soft,arc=2mm,left=3mm,right=3mm,top=1.5mm,bottom=1.5mm]
\small Exercicis graduats per nivells de dificultat:\quad \estrelles{1}\ bàsic \qquad \estrelles{2}\ mitjà \qquad \estrelles{3}\ avançat.\\[2pt] Resol-los al quadern mostrant tots els passos. Les taules i les rectes es poden completar directament en la fitxa. Fes els càlculs sense calculadora.
\end{tcolorbox}

% =====================================================
\seccio{1}{Nombres decimals}

\ex{1} Escriu com es llegeixen aquests nombres:\begin{tasks}(3)
  \task $7,25$
  \task $0,408$
  \task $12,06$
  \task $3,5$
  \task $100,001$
  \task $0,09$
\end{tasks}\recorda{Llegim la part entera, la paraula \emph{unitats} i la part decimal amb el nom de l'\textbf{última} xifra: $3,045$ són 3 unitats i 45 mil·lèsimes.}

\ex{1} Escriu amb xifres:\begin{tasks}(2)
  \task Tres unitats i quaranta-cinc mil·lèsimes
  \task Dotze centèsimes
  \task Vint unitats i nou dècimes
  \task Set mil·lèsimes
  \task Cent unitats i cinc centèsimes
  \task Quatre unitats i tres-centes mil·lèsimes
\end{tasks}

\ex{1} Quant val la xifra indicada en cada nombre?\begin{tasks}(3)
  \task El 7 en $34,571$
  \task El 3 en $2,738$
  \task El 7 en $607,15$
  \task El 8 en $0,0809$
  \task El 4 en $15,4$
  \task El 6 en $9,206$
\end{tasks}

\ex{2} Completa:\begin{tasks}(2)
  \task $1$ unitat $=$ \rule{1cm}{0.4pt} dècimes
  \task $1$ dècima $=$ \rule{1cm}{0.4pt} centèsimes
  \task $3$ unitats $=$ \rule{1cm}{0.4pt} centèsimes
  \task $45$ centèsimes $=$ \rule{1cm}{0.4pt} unitats
  \task $250$ mil·lèsimes $=$ \rule{1cm}{0.4pt} dècimes
  \task $7$ dècimes $=$ \rule{1cm}{0.4pt} mil·lèsimes
\end{tasks}

\ex{2} Descompon cada nombre segons el valor de les seues xifres:\begin{tasks}(3)
  \task $45,312$
  \task $7,09$
  \task $120,506$
\end{tasks}

\ex{3} Escriu el nombre que té: a) 4 desenes, 6 dècimes i 2 mil·lèsimes; \ b) 8 centenes, 3 unitats i 5 centèsimes; \ c) 15 centèsimes i 3 mil·lèsimes. Quin d'aquests és igual a $0,15$: \ $0,150$, \ $0,015$ o \ $0,1500$?


% =====================================================
\seccio{2}{Comparar, representar i aproximar decimals}

\ex{1} Escriu $<$, $>$ o $=$:\begin{tasks}(3)
  \task $4,7\ \square\ 4,68$
  \task $0,09\ \square\ 0,1$
  \task $5,307\ \square\ 5,37$
  \task $12,5\ \square\ 12,50$
  \task $3,099\ \square\ 3,1$
  \task $0,45\ \square\ 0,405$
\end{tasks}\recorda{Per a comparar, igualem el nombre de decimals afegint zeros: $\ 3,5\ \text{?}\ 3,48\ \to\ 3,50>3,48$.}

\ex{1} Ordena de menor a major: \ $2,5;\ 2,05;\ 2,55;\ 2,505;\ 2,0;\ 2,15$

\ex{2} Escriu el nombre decimal que correspon a cada punt:\par\smallskip{\centering \begin{tikzpicture}[baseline=0]\draw[-latex,thick] (-0.3,0)--(12.50,0);\draw[thick] (0.000,0.20)--(0.000,-0.20);\node[below=2pt,font=\small] at (0.000,-0.1) {$3$};\draw[muted] (1.200,0.12)--(1.200,-0.12);\draw[muted] (2.400,0.12)--(2.400,-0.12);\draw[muted] (3.600,0.12)--(3.600,-0.12);\draw[muted] (4.800,0.12)--(4.800,-0.12);\draw[muted] (6.000,0.12)--(6.000,-0.12);\draw[muted] (7.200,0.12)--(7.200,-0.12);\draw[muted] (8.400,0.12)--(8.400,-0.12);\draw[muted] (9.600,0.12)--(9.600,-0.12);\draw[muted] (10.800,0.12)--(10.800,-0.12);\draw[thick] (12.000,0.20)--(12.000,-0.20);\node[below=2pt,font=\small] at (12.000,-0.1) {$4$};\fill[accent] (2.400,0) circle (0.09) node[above=3pt,text=accent,font=\small]{A};\fill[accent] (8.400,0) circle (0.09) node[above=3pt,text=accent,font=\small]{B};\fill[accent] (10.800,0) circle (0.09) node[above=3pt,text=accent,font=\small]{C};\end{tikzpicture}\par}\par\smallskip{\centering \begin{tikzpicture}[baseline=0]\draw[-latex,thick] (-0.3,0)--(12.50,0);\draw[thick] (0.000,0.20)--(0.000,-0.20);\node[below=2pt,font=\small] at (0.000,-0.1) {$0,5$};\draw[muted] (1.200,0.12)--(1.200,-0.12);\draw[muted] (2.400,0.12)--(2.400,-0.12);\draw[muted] (3.600,0.12)--(3.600,-0.12);\draw[muted] (4.800,0.12)--(4.800,-0.12);\draw[muted] (6.000,0.12)--(6.000,-0.12);\draw[muted] (7.200,0.12)--(7.200,-0.12);\draw[muted] (8.400,0.12)--(8.400,-0.12);\draw[muted] (9.600,0.12)--(9.600,-0.12);\draw[muted] (10.800,0.12)--(10.800,-0.12);\draw[thick] (12.000,0.20)--(12.000,-0.20);\node[below=2pt,font=\small] at (12.000,-0.1) {$0,6$};\fill[accent] (3.600,0) circle (0.09) node[above=3pt,text=accent,font=\small]{D};\fill[accent] (9.600,0) circle (0.09) node[above=3pt,text=accent,font=\small]{E};\end{tikzpicture}\par}

\ex{2} Representa en la recta $1,4$ \ i \ $1,75$ (la recta va de 1 a 2).\par\smallskip{\centering \begin{tikzpicture}[baseline=0]\draw[-latex,thick] (-0.3,0)--(12.50,0);\draw[thick] (0.000,0.20)--(0.000,-0.20);\node[below=2pt,font=\small] at (0.000,-0.1) {$1$};\draw[muted] (0.600,0.12)--(0.600,-0.12);\draw[muted] (1.200,0.12)--(1.200,-0.12);\draw[muted] (1.800,0.12)--(1.800,-0.12);\draw[muted] (2.400,0.12)--(2.400,-0.12);\draw[muted] (3.000,0.12)--(3.000,-0.12);\draw[muted] (3.600,0.12)--(3.600,-0.12);\draw[muted] (4.200,0.12)--(4.200,-0.12);\draw[muted] (4.800,0.12)--(4.800,-0.12);\draw[muted] (5.400,0.12)--(5.400,-0.12);\draw[muted] (6.000,0.12)--(6.000,-0.12);\draw[muted] (6.600,0.12)--(6.600,-0.12);\draw[muted] (7.200,0.12)--(7.200,-0.12);\draw[muted] (7.800,0.12)--(7.800,-0.12);\draw[muted] (8.400,0.12)--(8.400,-0.12);\draw[muted] (9.000,0.12)--(9.000,-0.12);\draw[muted] (9.600,0.12)--(9.600,-0.12);\draw[muted] (10.200,0.12)--(10.200,-0.12);\draw[muted] (10.800,0.12)--(10.800,-0.12);\draw[muted] (11.400,0.12)--(11.400,-0.12);\draw[thick] (12.000,0.20)--(12.000,-0.20);\node[below=2pt,font=\small] at (12.000,-0.1) {$2$};\end{tikzpicture}\par}

\ex{2} Escriu tres nombres decimals compresos entre:\begin{tasks}(3)
  \task $0,3$ i $0,4$
  \task $2,71$ i $2,72$
  \task $5$ i $5,1$
\end{tasks}

\ex{2} Completa la taula aproximant cada nombre:\par\smallskip{\centering\renewcommand{\arraystretch}{1.35}\begin{tabular}{|c|C{2.3cm}|C{2.3cm}|C{2.3cm}|}\hline Nombre & Truncat a les centèsimes & Arrodonit a les dècimes & Arrodonit a les centèsimes\\\hline $3,4682$ & & & \\\hline $15,8962$ & & & \\\hline $0,7049$ & & & \\\hline $29,996$ & & & \\\hline\end{tabular}\par}\recorda{Per a \textbf{arrodonir}, mirem la primera xifra que eliminem: si és $0,1,2,3,4$ la deixem igual; si és $5,6,7,8,9$ sumem 1 a l'última xifra que queda.}

\ex{3} Un article costa $12,4583$\,€ segons la calculadora. Quant pagarem? Escriu un nombre que, arrodonit a les dècimes, siga $2,5$ i que siga menor que $2,5$. Quin és el nombre més xicotet amb dues xifres decimals que arrodonit a les unitats dona 8?


% =====================================================
\seccio{3}{Sumar, restar i multiplicar decimals}

\ex{1} Calcula:\begin{tasks}(3)
  \task $23,45+7,8$
  \task $40-12,35$
  \task $0,9+0,35$
  \task $15-6,47$
  \task $3,06+12,9+0,4$
  \task $100-0,01$
\end{tasks}\recorda{Per a sumar i restar, col·loquem les \textbf{comes alineades} i completem amb zeros.}

\ex{1} Calcula:\begin{tasks}(3)
  \task $4,6\cdot 3$
  \task $2,35\cdot 0,4$
  \task $0,6\cdot 0,05$
  \task $12,5\cdot 1,2$
  \task $3,25\cdot 1,4$
  \task $0,07\cdot 0,3$
\end{tasks}\recorda{Per a multiplicar, ho fem com si no hi haguera comes i separem tants decimals com \textbf{sumen} els dos factors.}

\ex{1} Calcula mentalment:\begin{tasks}(4)
  \task $3,7\cdot 100$
  \task $0,052\cdot 1000$
  \task $45,2:1000$
  \task $8:100$
  \task $6,15\cdot 10$
  \task $0,3:10$
  \task $250\cdot 0,1$
  \task $1,4\cdot 0,01$
\end{tasks}

\ex{2} Col·loca la coma en el resultat sense fer l'operació: \ $3,8\cdot2,5=950$; \ $0,45\cdot1,2=540$; \ $12,6\cdot0,03=378$. Comprova-ho estimant.

\ex{2} Completa amb el nombre que falta:\begin{tasks}(3)
  \task $2,75+\square=5$
  \task $\square-1,8=3,45$
  \task $0,6\cdot\square=60$
\end{tasks}

\ex{2} Comprem $2,5$ kg de pomes a $1,80$\,€/kg i $1,5$ kg de cireres a $6,40$\,€/kg. Paguem amb un bitllet de 20\,€. Quant ens tornen?

\ex{3} Una piscina fa $25$ m de llarg. La Laia fa $8$ llargs i el seu germà, $5,5$ llargs. Quants metres més ha nadat la Laia? Si la Laia nada cada llarg en $0,45$ minuts, quant tarda en total?

\ex{3} El perímetre d'un triangle és $17,25$ cm. Dos costats mesuren $4,8$ cm i $6,07$ cm. Quant mesura el tercer costat? Quant mesuraria el perímetre d'un triangle amb tots els costats 10 vegades més llargs?


% =====================================================
\seccio{4}{Dividir amb decimals}

\ex{1} Calcula:\begin{tasks}(3)
  \task $8,4:4$
  \task $9,75:3$
  \task $6,24:3$
  \task $0,8:2$
  \task $15,3:5$
  \task $1,26:6$
\end{tasks}

\ex{1} Calcula traient decimals al quocient fins que la divisió siga exacta:\begin{tasks}(3)
  \task $7:4$
  \task $15:8$
  \task $9:4$
  \task $3:5$
  \task $1:8$
  \task $21:12$
\end{tasks}\recorda{Si sobra residu, posem una coma al quocient, afegim un 0 al residu i continuem dividint.}

\ex{2} Calcula (si el divisor té decimals, multiplica dividend i divisor per 10, 100…):\begin{tasks}(3)
  \task $4,5:0,9$
  \task $2,73:0,7$
  \task $6:0,12$
  \task $12,6:0,35$
  \task $1,44:1,2$
  \task $0,9:0,03$
\end{tasks}

\ex{2} Calcula amb dues xifres decimals i indica el residu: \ a) $10:3$ \quad b) $5:7$ \quad c) $25,4:6$

\ex{2} Amb una botella de $1,5$ L, quants gots de $0,25$ L podem omplir? I amb una garrafa de 5 L?

\ex{3} Un tros de formatge de $0,85$ kg costa $11,39$\,€. Quin és el preu d'un quilo? Quant costaria un tros de $1,2$ kg?

\ex{3} Un cotxe ha gastat $31,5$ L de gasolina en un viatge de 450 km. Quants litres gasta cada 100 km? Si la gasolina costa $1,56$\,€/L, quant ha costat el viatge?


% =====================================================
\seccio{5}{Decimals i fraccions}

\ex{1} Escriu com a fracció decimal i després com a fracció irreductible:\begin{tasks}(4)
  \task $0,6$
  \task $0,45$
  \task $1,2$
  \task $0,004$
  \task $0,125$
  \task $2,75$
  \task $0,08$
  \task $3,5$
\end{tasks}\recorda{Al numerador posem el nombre sense coma; al denominador, un 1 seguit de tants zeros com decimals té: \ $0,35=\dfrac{35}{100}=\dfrac{7}{20}$.}

\ex{1} Escriu en forma decimal:\begin{tasks}(4)
  \task $\dfrac{7}{10}$
  \task $\dfrac{35}{100}$
  \task $\dfrac{4}{1000}$
  \task $\dfrac{125}{100}$
\end{tasks}

\ex{2} Escriu en forma decimal i classifica en exacte, periòdic pur o periòdic mixt:\begin{tasks}(4)
  \task $\dfrac{9}{20}$
  \task $\dfrac{5}{9}$
  \task $\dfrac{7}{12}$
  \task $\dfrac{3}{8}$
  \task $\dfrac{4}{11}$
  \task $\dfrac{1}{6}$
  \task $\dfrac{13}{5}$
  \task $\dfrac{2}{3}$
\end{tasks}\recorda{Dividim el numerador entre el denominador. Si les xifres es repeteixen, el decimal és \textbf{periòdic} i marquem el període: \ $\dfrac{5}{6}=0,8333\ldots=0,8\overline{3}$.}

\ex{2} Ordena de menor a major passant-ho tot a decimal: \ $0,7;\ \ \dfrac{2}{3};\ \ 0,67;\ \ \dfrac{3}{4};\ \ 0,6$

\ex{2} Quin és més gran?\begin{tasks}(3)
  \task $\dfrac{5}{8}$ o $0,62$
  \task $\dfrac{1}{3}$ o $0,33$
  \task $\dfrac{7}{5}$ o $1,4$
\end{tasks}

\ex{3} Escriu una fracció de denominador 3 compresa entre $1,5$ i $2$, i una fracció de denominador 8 compresa entre $0,3$ i $0,4$.

\ex{3} Un pastís es divideix en 8 porcions iguals. La Marta se'n menja $0,25$ del pastís i el Pau 3 porcions. Qui n'ha menjat més? Quina part del pastís queda, en fracció i en decimal?


% =====================================================
\seccio{6}{Raons i proporcions}

\ex{1} En una cistella hi ha 18 pomes i 24 peres. Calcula la raó entre pomes i peres i la raó entre peres i fruites totals. Què signifiquen?

\ex{1} Comprova si formen proporció:\begin{tasks}(4)
  \task $\dfrac{6}{9}$ i $\dfrac{8}{12}$
  \task $\dfrac{4}{10}$ i $\dfrac{6}{15}$
  \task $\dfrac{2,5}{4}$ i $\dfrac{5}{9}$
  \task $\dfrac{3}{7}$ i $\dfrac{12}{28}$
\end{tasks}\recorda{$\dfrac{a}{b}=\dfrac{c}{d}$ són proporció si $a\cdot d=b\cdot c$.}

\ex{2} Calcula el terme desconegut:\begin{tasks}(3)
  \task $\dfrac{3}{5}=\dfrac{12}{x}$
  \task $\dfrac{x}{8}=\dfrac{15}{24}$
  \task $\dfrac{4}{x}=\dfrac{10}{25}$
  \task $\dfrac{2,5}{4}=\dfrac{x}{12}$
  \task $\dfrac{1,2}{0,3}=\dfrac{8}{x}$
  \task $\dfrac{7}{x}=\dfrac{21}{6}$
\end{tasks}

\ex{2} Per a fer una massa de pizza fan falta 300 g de farina per cada 180 mL d'aigua. Quanta aigua necessitem per a 750 g de farina?

\ex{3} La raó entre l'edat d'una mare i la de la seua filla és $\dfrac{7}{2}$. Si la filla té 12 anys, quants en té la mare? I quants anys tenia la mare quan la filla en tenia 4? Es manté la raó?


% =====================================================
\seccio{7}{Proporcionalitat directa}

\ex{1} Completa la taula de proporcionalitat directa i calcula la constant de proporcionalitat:\par\smallskip{\centering {\renewcommand{\arraystretch}{1.35}\begin{tabular}{|l|C{1.25cm}|C{1.25cm}|C{1.25cm}|C{1.25cm}|C{1.25cm}|}\hline Entrades & 1 & 2 & 4 &  & 10\\\hline Preu (€) &  & 13 &  & 45,50 & \\\hline\end{tabular}}\par}

\ex{1} Indica si són directament proporcionals:\begin{tasks}(2)
  \task Les barres de pa que comprem i el preu que paguem.
  \task L'edat d'una persona i la seua altura.
  \task El costat d'un quadrat i el seu perímetre.
  \task Els litres de gasolina i els km recorreguts (amb el mateix cotxe).
\end{tasks}

\ex{1} Si 4 bolígrafs costen $3,20$\,€, quant costen 9 bolígrafs? I 15?

\ex{2} Un cotxe gasta $6,5$ L de gasolina cada 100 km. Quants litres gastarà en un viatge de 340 km?

\ex{2} Una aixeta omple 120 L en 8 minuts. Quant tardarà a omplir un dipòsit de 450 L? Quants litres haurà abocat en mitja hora?

\ex{2} Una recepta de pa de pessic per a 6 persones porta 250 g de farina, 4 ous i $0,15$ L de llet. Quines quantitats necessitem per a 9 persones?

\ex{3} En un mapa, 2 cm representen 15 km reals. Quina distància real hi ha entre dues ciutats que en el mapa estan a $9,4$ cm? A quants cm del mapa estan dos pobles separats per 6 km?

\ex{3} Un taxi cobra 3\,€ per pujar-hi i $1,20$\,€ per km. Completa la taula i digues si el preu és directament proporcional als km.\par\smallskip{\centering {\renewcommand{\arraystretch}{1.35}\begin{tabular}{|l|C{1.25cm}|C{1.25cm}|C{1.25cm}|C{1.25cm}|}\hline km & 1 & 2 & 5 & 10\\\hline Preu (€) &  &  &  & \\\hline\end{tabular}}\par}


% =====================================================
\seccio{8}{Proporcionalitat inversa}

\ex{1} Completa la taula de proporcionalitat inversa i calcula la constant:\par\smallskip{\centering {\renewcommand{\arraystretch}{1.35}\begin{tabular}{|l|C{1.25cm}|C{1.25cm}|C{1.25cm}|C{1.25cm}|C{1.25cm}|}\hline Velocitat (km/h) & 40 & 60 & 80 &  & 120\\\hline Temps (h) & 6 &  &  & 2,5 & \\\hline\end{tabular}}\par}

\ex{1} Indica si són directament proporcionals, inversament proporcionals o cap de les dues:\begin{tasks}(2)
  \task El nombre d'obrers i el temps que tarden a fer una paret.
  \task Els litres de pintura i els m$^2$ pintats.
  \task El nombre d'amics que paguen un regal i el que paga cadascun.
  \task El pes d'una persona i la seua edat.
\end{tasks}

\ex{1} Si 3 aixetes omplin una piscina en 10 hores, quant tardaran 5 aixetes iguals? I 2 aixetes?

\ex{2} Un grup de 12 persones té menjar per a 30 dies. Si se n'hi afegeixen 8 més, per a quants dies tindran menjar?

\ex{2} Un tren a 90 km/h tarda 2 hores a fer un trajecte. Quant tardaria a 120 km/h? A quina velocitat hauria d'anar per a tardar-ne 3?

\ex{3} 5 màquines fan una faena en 12 hores. Quantes màquines calen per a fer-la en 4 hores? Si només en tenim 3, quant tardaran?

\ex{3} Per a repartir l'aigua d'un dipòsit fan falta 60 garrafes de 5 L. Quantes garrafes de $1,5$ L farien falta? Si només tenim 25 garrafes iguals, de quina capacitat haurien de ser?


% =====================================================
\seccio{9}{Percentatges}

\ex{1} Completa la taula:\par\smallskip{\centering\renewcommand{\arraystretch}{1.5}\begin{tabular}{|c|C{2.4cm}|C{2.4cm}|}\hline Percentatge & Fracció irreductible & Decimal\\\hline $40\,\%$ & & \\\hline $5\,\%$ & & \\\hline $75\,\%$ & & \\\hline $12\,\%$ & & \\\hline $120\,\%$ & & \\\hline $2,5\,\%$ & & \\\hline\end{tabular}\par}

\ex{1} Calcula:\begin{tasks}(3)
  \task El $20\,\%$ de $45$
  \task El $75\,\%$ de $120$
  \task El $15\,\%$ de $60$
  \task El $8\,\%$ de $350$
  \task El $12\,\%$ de $250$
  \task El $7,5\,\%$ de $60$
\end{tasks}\recorda{El $p\,\%$ de $N$ és $N\cdot\dfrac{p}{100}$. Per exemple, el $35\,\%$ de $80$ és $80\cdot0,35=28$.}

\ex{1} Expressa en percentatge:\begin{tasks}(6)
  \task $0,35$
  \task $0,06$
  \task $1,5$
  \task $\dfrac{3}{5}$
  \task $\dfrac{7}{20}$
  \task $\dfrac{1}{8}$
\end{tasks}

\ex{2} En un institut de 650 alumnes, 91 vénen amb bicicleta. Quin percentatge representen? I en una classe de 25 alumnes, si 6 porten ulleres?

\ex{2} En una classe de 28 alumnes, el $25\,\%$ ha tret excel·lent i el $50\,\%$ notable. Quants alumnes han tret cada nota? Quin percentatge no ha tret ni excel·lent ni notable?

\ex{3} Hem llegit 135 pàgines d'un llibre, que són el $45\,\%$ del total. Quantes pàgines té el llibre?

\ex{3} Què és més, el $30\,\%$ de 50 o el $50\,\%$ de 30? I el $20\,\%$ de 80 o el $80\,\%$ de 20? Explica per què passa.


% =====================================================
\seccio{10}{Augments i disminucions percentuals}

\ex{1} Escriu l'índex de variació:\begin{tasks}(3)
  \task augment del 30 \%
  \task disminució del 8 \%
  \task augment del 4 \%
  \task rebaixa del 50 \%
  \task IVA del 21 \%
  \task descompte del 15 \%
\end{tasks}\recorda{Augment del $p\,\%$: multipliquem per $1+\dfrac{p}{100}$. \ Disminució del $p\,\%$: multipliquem per $1-\dfrac{p}{100}$.}

\ex{1} Calcula la quantitat final:\begin{tasks}(2)
  \task $80$\,€ amb un descompte del $15\,\%$
  \task $50$\,€ amb una pujada del $21\,\%$
  \task $650$\,€ amb un descompte del $12\,\%$
  \task $18$\,€ amb una pujada del $4\,\%$
  \task $2400$\,€ amb una pujada del $5\,\%$
  \task $64,90$\,€ amb un descompte del $30\,\%$
\end{tasks}

\ex{2} Una samarreta passa de 25\,€ a 20\,€. Quin percentatge de descompte li han fet? I si un abric passa de 120\,€ a 138\,€, quin percentatge ha pujat?

\ex{2} En una botiga, un televisor costa 450\,€ sense IVA. Quant costa amb l'IVA del 21\,\%? Si a més ens fan un descompte del 10\,\% sobre el preu amb IVA, quant paguem?

\ex{3} Després d'aplicar l'IVA del 21\,\%, una factura és de $108,90$\,€. Quant valia sense IVA? I si després d'una rebaixa del 15\,\% paguem 68\,€, quin era el preu inicial?

\ex{3} El preu d'un patinet de 200\,€ puja un 10\,\% i, al cap d'un mes, baixa un 10\,\%. Quant costa al final? Torna al preu inicial? Què passaria si primer baixara un 20\,\% i després pujara un 20\,\%?

\ex{3} La població d'un poble era de 2400 habitants. El primer any augmenta un 5\,\% i el segon any disminueix un 5\,\%. Quants habitants té al final? Quin percentatge ha variat respecte a l'inici?

\vfill
\begin{center}\color{muted}\small\faCheckCircle[regular]\ \ Ànim! Si et costa algun exercici, repassa els apunts interactius de la unitat.\end{center}

\end{document}
